\chapter{Analytical Idealism and the Ontology of Mind}
\label{ch:analytical-idealism}

\epigraph{\itshape The whole of nature, including ourselves and our
inner experiences, is the extrinsic appearance of a unitary mental
process.}{--- Bernardo Kastrup}

\noindent
This chapter develops the third major resource of Part II:
contemporary analytical idealism, particularly in the form
articulated by the Dutch philosopher Bernardo Kastrup. The treatment
serves two purposes within the larger project. First, it provides
the contemporary philosophical vocabulary in which the Akbarian
metaphysical claims of Chapter 9 can be expressed and compared with
modern Western philosophy of mind. Second, it removes the
materialist obstacle that has made Islamic metaphysical claims about
provision, blessing, intention, and the nature of wealth appear
implausible to modern academic readers.

The chapter assumes no background in philosophy of mind. It develops
the central problems --- particularly the so-called ``hard problem
of consciousness'' --- from the ground up, surveys the major
positions, and then concentrates on idealism in its analytical form.
The chapter is more philosophical than mathematical; the reader
should expect arguments rather than equations. Where the philosophy
admits formal expression we provide it; where it does not, we are
careful to be clear about what is being claimed.

The chapter proceeds in nine sections.
\xref{sec:hard-problem} presents the hard problem of consciousness.
\xref{sec:positions-overview} surveys the major positions in
contemporary philosophy of mind.
\xref{sec:materialism-difficulties} examines the difficulties of
materialism (the dominant academic position).
\xref{sec:dualism-panpsychism} treats property dualism and
panpsychism as alternatives.
\xref{sec:analytical-idealism-core} presents Kastrup's analytical
idealism in detail.
\xref{sec:dissociation} treats the dissociation analogy from clinical
psychiatry that anchors the analytical-idealist account.
\xref{sec:empirical-anomalies} surveys empirical anomalies that
challenge materialism.
\xref{sec:tawhid-correspondence} draws the structural correspondence
with Islamic \arb{tawhid}.
\xref{sec:idealism-ahead} previews how the Akbarian framework of
Chapter 9 will extend the contemporary idealist position.

\section{The hard problem of consciousness}
\label{sec:hard-problem}

\subsection{Statement of the problem}

In 1995 the Australian philosopher David Chalmers introduced what is
now the standard formulation of the central puzzle in philosophy of
mind:\footnote{Chalmers, D. (1995), ``Facing up to the problem of
consciousness,'' \emph{Journal of Consciousness Studies} 2(3):
200--219.}

\begin{conceptbox}[The hard problem of consciousness]
\label{conc:hard-problem}
A complete physical description of the brain ---
every neuron, every synapse, every electrochemical signal --- does
not appear to entail the existence of subjective experience. The
neuroscientist can describe how visual information enters the
retina, how it is processed by the visual cortex, how it is
associated with motor responses; none of this description \emph{is}
the experience of seeing red. The subjective character of
experience --- what philosophers call \emph{qualia} --- is
unaccounted for in the physical description. Why is there
something it is like to be a brain?
\end{conceptbox}

Chalmers's distinction between the ``easy problems'' (explaining
specific cognitive functions in terms of brain mechanisms) and the
``hard problem'' (explaining why any of this gives rise to
subjective experience at all) has structured the debate for the
past three decades. The easy problems are, in principle, soluble by
neuroscience: explain attention, memory, perception, motor control,
language by identifying the neural mechanisms. The hard problem is
not in principle soluble by such methods, because it asks not
\emph{how} the brain does something but \emph{why} doing it produces
subjective experience.

\subsection{Why the problem is genuinely hard}

The depth of the hard problem becomes clear when we consider what a
solution would have to look like. A solution would have to specify a
physical state $S$ and a subjective experience $E$, and explain why
$S$ \emph{must} be accompanied by $E$. The explanation cannot be
``$S$ is correlated with $E$'' (we already know that brain states
correlate with experiences); the explanation must be why the
correlation must hold, why a being in state $S$ must experience $E$,
why a being in state $S$ could not exist without experiencing $E$.

No such explanation has been forthcoming. The standard physicalist
response is to claim that the explanation will eventually be found:
the problem is hard because we don't yet have it, but science will
solve it as it has solved other hard problems. This response is
sometimes called ``promissory materialism.'' It is unfalsifiable in
its form (any failure to solve the problem can be attributed to
incomplete science), and it has been waiting for redemption for at
least three decades without any sign that it is coming.

A more honest response is to acknowledge that the problem may be
genuinely intractable for any physicalist framework. If subjective
experience is not entailed by physical description, then no
extension of physical description --- however detailed, however
mathematically sophisticated --- will produce it. The conclusion may
be that something other than physical description is required.

\subsection{The zombie argument}

Chalmers and others have made the case for the genuine intractability
through the so-called \emph{zombie argument}. A philosophical
zombie is a being physically and behaviorally identical to a normal
human but lacking subjective experience. The being walks, talks,
responds to questions, and has all the same brain processes as a
normal human; but there is nothing it is like to be the being.

The argument: such zombies are at least conceivable (we can imagine
them without contradiction). If conceivability is a guide to
possibility, then zombies are possible. If zombies are possible,
then subjective experience is not entailed by physical description
(since zombies have the physical description but lack the
experience). If subjective experience is not entailed by physical
description, then physicalism is false.

Each step in the argument has been disputed. Some philosophers hold
that zombies are not conceivable on close examination. Others hold
that conceivability is not a reliable guide to possibility. Still
others hold that physicalism does not require subjective experience
to be entailed by physical description in the strong sense the
argument assumes.

For our purposes, we need not adjudicate the zombie argument. What
matters is that the argument is taken seriously by working
philosophers and that no consensus solution has emerged. The hard
problem is, by professional consensus, hard.

\section{Major positions in contemporary philosophy of mind}
\label{sec:positions-overview}

The contemporary landscape contains several positions on the
mind-body problem. We survey them briefly to situate analytical
idealism in context.

\subsection{Materialism (physicalism)}

The dominant position in twentieth-century academic philosophy.
Materialism holds that everything that exists is physical (or
supervenient on the physical). Mental states, including conscious
experiences, are identical to or constituted by physical states of
the brain.

Within materialism, a major distinction is between
\emph{type-identity theory} (mental state types are identical to
physical state types: pain just is C-fiber firing) and
\emph{functionalism} (mental states are characterized by their
functional roles; whatever physical state plays the right role
counts as the mental state).

Functionalism is currently more popular than type-identity theory.
Both face the hard problem: they describe what mental states do but
not why they have the subjective character they do.

\subsection{Property dualism}

Property dualism holds that there is one substance (the physical)
but two kinds of properties (physical and mental). The mental
properties are not reducible to the physical, even though they are
properties of physical things (typically brains).

Chalmers himself defends a version of property dualism. The position
is internally consistent but faces its own difficulties: how do
mental and physical properties interact, given that they are
distinct? Why do mental properties supervene on physical properties
in the patterns they do?

\subsection{Panpsychism}

Panpsychism is the view that consciousness is a fundamental feature
of matter, present in some form in all things, however attenuated.
Particles have proto-experiences; brains, as complex aggregates of
particles, have complex experiences.

Panpsychism has gained significant attention in recent decades, with
serious defenses by Galen Strawson (who argues that physicalism
collapses into panpsychism when consistently held), Philip Goff,
and others. Chalmers's later work is sympathetic to panpsychism.

The major difficulty is the \emph{combination problem}: how do the
elementary consciousnesses of fundamental particles combine into
the unified consciousness of a brain? The unity of conscious
experience appears to require some additional principle that
panpsychism does not naturally supply.

\subsection{Idealism}

Idealism holds that mind is fundamental and matter is derivative.
Physical reality is, in some way, an appearance within or to mind;
mental reality does not need to be explained by reference to
physical reality, because it is the more basic ontological category.

Idealism was the dominant position in much of nineteenth-century
philosophy (Hegel, Schelling, Bradley, Royce, Berkeley earlier).
It was largely abandoned in twentieth-century Anglo-American
philosophy, partly because of perceived associations with religion
and partly because of the rise of physicalism.

The position is now returning, both in panpsychist forms (which can
be read as a version of idealism) and in more direct forms,
particularly in the analytical idealism of Bernardo Kastrup. We
develop Kastrup's version in detail in \xref{sec:analytical-idealism-core}.

\subsection{Other positions}

For completeness, we mention briefly:

\begin{itemize}[leftmargin=2em]
    \item \emph{Eliminativism}: the view that mental states do not
    really exist and that folk-psychological vocabulary will be
    replaced by neuroscientific vocabulary. Defended by the
    Churchlands. Few accept it.
    \item \emph{Mysterianism}: the view that consciousness is a real
    phenomenon but our cognitive faculties are unable to understand
    it. Defended by Colin McGinn. Concedes the hard problem and
    denies its solubility.
    \item \emph{Neutral monism}: the view that mind and matter are
    aspects of a more basic substance that is neither. Defended by
    Russell, Mach, and James. Revived in contemporary form by some
    panpsychists.
\end{itemize}

The reader who wants a comprehensive survey is referred to Chalmers
(2010) \emph{The Character of Consciousness}.

\section{The difficulties of materialism}
\label{sec:materialism-difficulties}

We focus on materialism because it is the dominant position and
because the case against analytical idealism is largely the case for
materialism. The case for materialism, however, is weaker than its
dominance suggests. We summarize the major difficulties.

\subsection{The hard problem itself}

We have already discussed the hard problem (\xref{sec:hard-problem}).
It is the central difficulty for materialism. Materialism holds
that subjective experience is somehow produced by physical processes;
no specification of how has been forthcoming; the absence has
persisted for centuries; the philosophical pressure is increasing
rather than decreasing.

\subsection{The unity of consciousness}

Subjective experience appears to be unified: at any moment, my
visual experience, auditory experience, bodily experience, emotional
experience, and ongoing thoughts are integrated into a single field
of awareness. This unity is not a small matter; it is the structural
character of conscious experience.

Materialism, particularly in its modern functionalist forms, has
difficulty explaining the unity. The brain is composed of many
different processing modules operating in parallel; how they combine
into a unified experience is unclear. The ``binding problem'' in
neuroscience is the technical formulation of this difficulty;
proposed solutions (synchronous neural firing, global workspace
theories, integrated information theory) are partial and contested.

\subsection{Mental causation}

If mental states are identical to or supervenient on physical
states, what causal work do they do? It is tempting to say that
mental states cause behaviors --- I see the apple, I want it, I
reach for it --- but if the underlying physical states are causally
sufficient for the behaviors, the mental states are causally
redundant. This is the so-called \emph{exclusion problem}: physical
causes exclude non-physical (or supervenient) ones.

Various solutions have been proposed (compatibilist accounts of
mental causation, non-reductive physicalism, ``causal exclusion''
acceptance). None is fully satisfying. The intuitive view that
mental states do real causal work (we act because of our beliefs and
desires, not despite them) sits uneasily with the materialist
framework.

\subsection{Brain damage and cognitive function}

Materialism predicts that brain damage should correlate strongly
with cognitive impairment. The prediction is correct in many cases
but breaks down in others. Hydrocephalic patients with most of
their brain volume replaced by cerebrospinal fluid have been
documented with normal IQ.\footnote{Lewin, R. (1980), ``Is your
brain really necessary?,'' \emph{Science} 210: 1232--1234.
The case studies described by John Lorber at the University of
Sheffield are striking; one mathematics graduate student with an
IQ of 126 was found by CT scan to have ``virtually no brain''
(approximately 5\% of normal brain volume).} Patients who have lost
substantial cortical tissue can recover function more completely
than the materialist hypothesis predicts.

These cases do not refute materialism; they can be explained by
neural plasticity, distributed processing, and the brain's capacity
for compensation. But they sit uneasily with the strong-materialist
intuition that mind is just the brain's activity.

\subsection{Psychedelic experience and brain activity}

A more recent line of evidence: psychedelic experience.
Functional MRI studies of subjects under psilocybin and other
psychedelics consistently find that the experience of expanded,
intensified, more vivid consciousness corresponds to
\emph{decreased} brain activity, particularly in the so-called
default mode network.\footnote{Carhart-Harris, R. L. et al.\ (2012),
``Neural correlates of the psychedelic state as determined by fMRI
studies with psilocybin,'' \emph{PNAS} 109(6): 2138--2143.}

The materialist prediction would be the opposite: more vivid
experience requires more neural activity. The empirical finding ---
robust across multiple studies and multiple substances --- is that
the relationship is reversed. Defenders of materialism have proposed
explanations (the default mode network normally inhibits experience;
disinhibition produces expanded experience), but the explanations
require ad hoc structure that the materialist framework does not
naturally supply.

\subsection{Quantum-mechanical observer dependence}

We treated the quantum interpretive disputes in Chapter 5
(\xref{ch:quantum-foundations}). Several major interpretations
(QBism, certain forms of objective collapse) attribute a constitutive
role to the observer in the production of measurement outcomes.
This is at least awkward for materialism, which would prefer that
physical reality be observer-independent. The issue is not decisive
--- materialist interpretations of quantum mechanics exist --- but
it is part of the broader pattern: the empirical and theoretical
pressures on materialism are accumulating.

\subsection{Cumulative weight}

None of these difficulties is decisive in isolation. The
materialist can offer rejoinders to each. The cumulative weight,
however, is substantial. The hard problem is not solved. The unity
of consciousness is not explained. Mental causation is awkward.
Brain damage cases are anomalous. Psychedelic findings reverse the
prediction. Quantum mechanics adds awkwardness. None of this proves
materialism false; together they do show that the dominance of
materialism in twentieth-century academic philosophy was
sociological as much as epistemic.

\section{Property dualism and panpsychism: alternatives}
\label{sec:dualism-panpsychism}

Before turning to analytical idealism, we briefly examine the two
positions that have been the most prominent contemporary
alternatives to materialism.

\subsection{Property dualism}

Property dualism inherits the hard problem in modified form. If
mental and physical are distinct properties of the same substance,
then the question becomes: why does the physical substance have the
specific mental properties it does? Why does brain state $B$ have
the qualia $Q$ rather than some other set of qualia, or none at all?

Chalmers, who defends property dualism in some moods, calls these
``psychophysical bridge laws'' --- laws connecting physical states
to mental properties. The laws are taken as fundamental, not
derivable from anything more basic. The position is consistent but
has the cost of multiplying fundamental laws (we have physical laws
\emph{and} psychophysical bridge laws), with no explanation of why
the bridge laws take the form they do.

\subsection{Panpsychism and the combination problem}

Panpsychism faces the combination problem most acutely. If
consciousness is a feature of all matter, then a brain is just an
aggregate of conscious particles. But the brain has unified conscious
experience; the particles do not (or, if they do, their experiences
are minimal and unrelated to mine). How do the elementary
consciousnesses combine into the unified consciousness?

The combination problem is not a small technical issue. It is the
central difficulty of panpsychism, and proposed solutions
(``constitutive panpsychism,'' ``cosmopsychism,'' various
mereological proposals) are highly contested. The most influential
contemporary defender of panpsychism, Philip Goff, has written
extensively about the combination problem and acknowledges it as
the central challenge.\footnote{Goff, P. (2017), \emph{Consciousness
and Fundamental Reality}, Oxford University Press.}

\subsection{Why analytical idealism avoids these problems}

Idealism, in Kastrup's analytical form, avoids the hard problem by
denying that consciousness needs to arise from anything; consciousness
is fundamental. It avoids the combination problem by holding that
consciousness is unified at the most fundamental level (Mind-at-Large)
and that individual minds are dissociated portions of this unity,
not aggregates of smaller consciousnesses. It avoids the bridge-law
problem by treating physical reality as the appearance of mental
processes; the relationship is constitutive, not law-governed.

The advantages do not, of course, make idealism true. They show that
idealism is a serious option that handles certain difficulties
better than the alternatives. The case for idealism is positive as
well as negative; we develop it next.

\section{Analytical idealism: the core thesis}
\label{sec:analytical-idealism-core}

\subsection{Statement of the thesis}

\begin{conceptbox}[Analytical idealism]
\label{conc:analytical-idealism}
The fundamental ontological category is mind. There is a single,
universal mind, which Kastrup calls \emph{Mind-at-Large}. This mind
is the only existent in the basic sense.

Individual minds (yours, mine, those of other organisms) are
\emph{dissociated alters} of Mind-at-Large: localized portions of
the universal mind that have become functionally separated from the
rest, analogous to dissociated personalities in dissociative
identity disorder.

Physical reality --- including brains, bodies, and the rest of
matter --- is how Mind-at-Large appears \emph{extrinsically} to its
own dissociated alters. From the outside (from the perspective of
one alter looking at another, or at the rest of Mind-at-Large), the
mental processes of Mind-at-Large appear as physical structures and
processes. From the inside (from the perspective of the mental
process itself), they are experiences.

Physical reality and mental reality are not two distinct things;
they are two perspectives on the same thing. The physical is the
extrinsic appearance of the mental.
\end{conceptbox}

This thesis is a substantial revision of the modern Western default.
We unpack its components.

\subsection{Mind as fundamental}

The thesis that mind is the fundamental ontological category is the
inverse of materialism. Materialism takes the physical as
fundamental and asks how mind arises from it. Idealism takes the
mental as fundamental and asks how the physical (as appearance)
arises from it. Both are coherent positions; the question is which
is better supported by argument and evidence.

The argument for taking the mental as fundamental rests on the
hard problem: if mental experience cannot be derived from physical
description, then the mental is either fundamental or epiphenomenal;
the epiphenomenal options have their own difficulties (mental
causation, the unity of consciousness), so the fundamental option is
the cleaner choice.

A second argument: the only reality we have direct access to is
mental. We do not perceive the world directly; we have experiences,
which we then interpret as being caused by an external world. The
external world is an inference, not a direct given. The mental is
direct; the physical is inferred. To take the inferred as
fundamental and the direct as derivative is, on this argument,
backwards.

\subsection{Mind-at-Large}

The thesis is not that each individual mind is fundamental; that
would be a kind of solipsism (or, in its multi-mind version, would
face severe coordination problems among independent fundamental
minds). The thesis is that there is a single universal mind, and
that individual minds are aspects or portions of this single mind.

The terminology ``Mind-at-Large'' is Kastrup's; the same idea has
appeared in many philosophical traditions under different names ---
the One of Plotinus, the World Soul of Plato, the Absolute of
Hegel, the Brahman of Vedanta, the Tao of Daoism. The contemporary
form is distinguished by being articulated in modern philosophical
vocabulary and by attempting to do explicit analytical work.

In the Akbarian framework that we treat in Chapter 9, the
correspondent of Mind-at-Large is absolute \arb{wujud}: the
foundational ontological category of being, of which all created
things are loci of self-disclosure. The structural correspondence is
not coincidence; it reflects the fact that both traditions are
attempting to articulate the same insight. We develop this in
\xref{sec:tawhid-correspondence}.

\subsection{The dissociation analogy}

The most distinctive feature of Kastrup's analytical idealism is the
dissociation analogy. Individual minds are not parts of Mind-at-Large
in the way that physical objects can be parts of a whole; they are
\emph{dissociated alters} of Mind-at-Large. The analogy is to
dissociative identity disorder (DID), the clinical condition in
which a single human person manifests multiple distinct
personalities, often mutually unaware, each with its own
subjective experience.

The dissociation analogy is more than rhetorical. The clinical
literature on DID provides the empirical anchor for the
analytical-idealist account of how a single mind can fragment into
multiple mutually-isolated mental processes. We develop this in
\xref{sec:dissociation}.

\subsection{Physical reality as extrinsic appearance}

The thesis that physical reality is the extrinsic appearance of
mental processes is the most counterintuitive aspect of analytical
idealism. To say that the brain is what mental processes ``look
like'' from outside is to invert the usual intuition (that mental
processes are what brains \emph{do} from inside).

The inversion is supported by the following observation: from the
inside, mental processes feel like experiences (qualia, thoughts,
volitions). From the outside, the same processes appear as physical
structures (neurons firing, neurotransmitters flowing, behavior
manifesting). The two are not separate things connected by
mysterious laws; they are the same thing seen from two perspectives.

This is a strong claim. It requires that all physical reality ---
including reality far from any obvious mental process, like
distant galaxies or empty space --- be the extrinsic appearance of
some mental process within Mind-at-Large. Kastrup accepts this
implication. The mental processes of Mind-at-Large are vast and
complex; the physical reality we observe is what these processes
look like from the perspective of the dissociated alters that we
are.

\subsection{The empirical predictions of analytical idealism}

Analytical idealism makes several empirical predictions:

\begin{itemize}[leftmargin=2em]
    \item Brain processes correlate with experience but do not
    \emph{produce} experience. The correlation is between two
    perspectives on the same thing.
    \item The structure of physical reality should track the
    structure of the underlying mental process. This is consistent
    with the deep mathematical regularities of physical law.
    \item Anomalies in brain-behavior correlation (hydrocephalic
    cases, psychedelic findings) are predicted: the brain is the
    extrinsic appearance, but the underlying mental process is what
    is doing the work.
    \item Quantum-mechanical observer effects are natural: the
    observer is part of the underlying mental process, and the
    observation is the dissociated alter's interaction with another
    portion of Mind-at-Large.
\end{itemize}

These predictions are not decisive. Materialism can offer competing
predictions; competing interpretations of the data exist. The
predictions show that analytical idealism is empirically engaged,
not metaphysically idle.

\section{The dissociation analogy}
\label{sec:dissociation}

The dissociation analogy is the empirical anchor for Kastrup's
account. We treat it in some detail because it is doing significant
work.

\subsection{Dissociative identity disorder}

DID, formerly called multiple personality disorder, is a recognized
psychiatric condition in which a single person manifests multiple
distinct personalities (``alters'') that may be mutually unaware. The
condition is rare and contested in some quarters, but the clinical
literature is substantial and the diagnostic criteria are now
established.\footnote{For a contemporary overview, see Brand, B. L.
et al.\ (2016), ``Separating fact from fiction: an empirical
examination of six myths about dissociative identity disorder,''
\emph{Harvard Review of Psychiatry} 24(4): 257--270.}

Several features of DID are particularly relevant to the
analytical-idealist account.

First, the alters can have distinct subjective experiences. Each
alter has its own first-person perspective, often with its own
memories, emotional patterns, and characteristic responses. The
experiential isolation is not just behavioral; the alters report
distinct phenomenologies.

Second, the alters can have distinct physiological correlates.
Documented cases include alters with different visual acuity (one
near-sighted, another not), different allergies, different
handedness, and different responses to medication. The
physiological differences are within a single human body, but the
body's properties shift as different alters take executive
control.\footnote{Strasburger, H. \& Waldvogel, B. (2015),
``Sight and blindness in the same person: gating in the visual
system,'' \emph{PsyCh Journal} 4(4): 178--185.}

Third, the dissociation is reversible in some cases through
therapy. Alters can be integrated back into a single personality;
the mutual isolation can be overcome. This shows that the
dissociation is a functional separation, not a metaphysical
multiplication of selves.

\subsection{The analogy to Mind-at-Large}

Kastrup's analogy: individual minds are to Mind-at-Large as alters
are to the dissociated person. There is one underlying mind. It
appears, from the inside, as multiple mutually-isolated minds. The
isolation is functional, not metaphysical. Each individual mind has
its own first-person perspective, its own memories, its own
experiences --- but they are aspects of a single underlying
mental process.

The analogy gives the analytical-idealist account a specific
mechanism for how the unity of Mind-at-Large can be consistent with
the apparent multiplicity of individual minds. We are not separate
fundamental minds; we are dissociated portions of a single mind.

The analogy also gives the account empirical purchase. DID is a
real, documented phenomenon. The mechanisms by which a single mind
can fragment are studied clinically. The analytical-idealist
hypothesis extends these mechanisms to the cosmic scale, but the
extension is grounded in the well-attested smaller-scale phenomenon.

\subsection{Limitations of the analogy}

Kastrup is careful to note the analogy is not perfect. DID arises
under specific conditions (typically severe early trauma) and has
specific characteristics (the alters often have personality
patterns related to the trauma); these specifics do not transfer
cleanly to the cosmic scale. The analogy works at the structural
level (a single mind appears as multiple), not at the
detailed level.

The structural level is what matters for the analytical-idealist
argument. The argument is that a single mind \emph{can} appear as
multiple mutually-isolated mental processes; DID demonstrates this
empirically; the analytical-idealist hypothesis is that the
empirical demonstration is reproducible at the cosmic scale.

\section{Empirical anomalies for materialism}
\label{sec:empirical-anomalies}

We treat briefly several lines of empirical evidence that are
difficult for materialism but natural for analytical idealism.

\subsection{Brain damage and consciousness}

We have already noted the hydrocephalic cases (\xref{sec:materialism-difficulties}).
The materialist hypothesis predicts that severe brain damage should
produce severe cognitive impairment. The prediction is often correct
but sometimes breaks down dramatically. Patients with
hydranencephaly (essentially no cerebral cortex) have been
documented with rich experiential lives.\footnote{Merker, B.
(2007), ``Consciousness without a cerebral cortex,'' \emph{Behavioral
and Brain Sciences} 30(1): 63--81.}

The analytical-idealist explanation: the brain is the extrinsic
appearance of certain mental processes, but the underlying mental
processes are not exhausted by the cortical structure. When the
cortex is missing, the underlying mental processes can still be
present, manifesting through whatever extrinsic structure is
available.

\subsection{Near-death and out-of-body experiences}

The literature on near-death experiences (NDEs) is large and
contested. The phenomenon is well-attested: patients reporting
extended, vivid experiences during periods of cardiac arrest or
severe physical compromise, often including accurate perception of
events that the physical perspective should not have allowed. The
materialist framework has difficulty explaining how such
experiences can occur during periods when normal brain function is
absent.

The analytical-idealist framework provides a natural account: the
brain is the extrinsic appearance of certain mental processes, but
the mental processes themselves are not produced by the brain.
During cardiac arrest, the extrinsic appearance is disrupted; the
underlying mental processes continue, possibly in modes not normally
accessible.

The reader should take the NDE literature with appropriate caution.
Many of the cases are anecdotal; some are contested; alternative
explanations exist. The point is not that NDEs prove anything but
that they are anomalous for materialism and accommodated naturally
by analytical idealism.

\subsection{Psychedelic studies}

The psychedelic-fMRI findings (\xref{sec:materialism-difficulties})
are particularly striking. Subjects under psilocybin report
expanded, intensified experience; the brain shows decreased activity
in the default mode network. The materialist prediction (more
experience requires more neural activity) is reversed.

The analytical-idealist explanation: the brain (specifically, the
default mode network) normally constrains and filters the mental
process; reducing the constraint allows more of the underlying
mental process to manifest. The brain is a filter for consciousness,
not the producer of it.\footnote{This view, sometimes called the
``filter theory'' of consciousness, was articulated by William
James and Henri Bergson in the early twentieth century. Aldous
Huxley elaborated it in \emph{The Doors of Perception} (1954). It
has been revived in the contemporary psychedelic literature.}

\subsection{Quantum-mechanical observer dependence}

We treated this in Chapter 5. The key point: several major
interpretations of quantum mechanics give the observer a
constitutive role in the production of measurement outcomes. The
observer is not just measuring something that is independently
there; the observer participates in determining what is there. This
is awkward for materialism (which would prefer observer-independence)
and natural for idealism (which expects mind to be constitutive of
appearance).

\subsection{Cumulative weight}

As with the difficulties of materialism, none of these anomalies is
decisive. Each can be explained, with effort, within materialism.
Cumulatively, however, they suggest that the materialist framework
is being asked to do more work than it can comfortably do. Idealism
in its analytical form accommodates them more naturally.

\section{The structural correspondence with Islamic tawhid}
\label{sec:tawhid-correspondence}

We now draw the structural correspondence between Kastrup's
analytical idealism and the Islamic doctrine of \arb{tawhid}, which
will be developed in detail in Chapter 9.

\subsection{Tawhid as ontological monism}

The first principle of Islam --- \arb{la ilaha illa Allah}, ``there
is no being but Allah'' --- has been read in the Akbarian tradition
as a proposition not just about the exclusivity of worship but about
the exclusivity of being. In its strongest reading (the doctrine of
\arb{wahdat al-wujud} that we treat in Chapter 9), the proposition
asserts that there is, ultimately, only one being; created things
are not independent existents but loci of divine self-disclosure.

This is, structurally, the same claim as Kastrup's analytical
idealism. There is one fundamental ontological reality
(Mind-at-Large for Kastrup, absolute \arb{wujud} for Ibn Arabi);
all apparent multiplicity is the fragmentation of this unity into
mutually-isolated portions (dissociated alters for Kastrup, loci
of self-disclosure for Ibn Arabi).

The correspondence is not coincidence. It is what we should expect
when two traditions, working from different starting points,
attempt to articulate the same underlying truth. The Islamic
tradition, working from revealed texts and contemplative practice,
arrived at the doctrine in the seventh through thirteenth Islamic
centuries. The contemporary Western philosophical tradition, working
from the hard problem of consciousness and the interpretive disputes
of quantum mechanics, has arrived at a structurally similar
position.

\subsection{Differences in theological content}

The convergence is structural, not theological. Kastrup's
Mind-at-Large is not the Islamic Allah. Specifically:

\begin{itemize}[leftmargin=2em]
    \item Kastrup makes no claim that Mind-at-Large is personal,
    benevolent, or the source of moral law. The Islamic tradition
    does. Mind-at-Large might or might not be all of these things;
    Kastrup's analysis is silent.
    \item Kastrup makes no claim about the relationship between
    Mind-at-Large and revelation, prophecy, or scripture. The
    Islamic tradition makes specific claims here.
    \item Kastrup makes no claim about the eschatological structure
    of the universe (the afterlife, judgment, and so on). The
    Islamic tradition does.
\end{itemize}

The differences are substantial. The convergence is in the basic
ontology: mind is fundamental; physical reality is appearance of
mental process; individual minds are not separate fundamental
realities but localized aspects of a unified underlying mind.

The Islamic tradition adds the further theological content. The
addition is not arbitrary --- it follows from the specific
revelatory framework of Islam --- but it goes beyond what Kastrup's
philosophical argument establishes.

\subsection{What this enables for the project}

The structural correspondence enables a key argumentative move that
the rest of this book will use repeatedly. We can show that the
Islamic metaphysical claims, when read through the Akbarian
framework, correspond to a coherent ontological position
(analytical idealism in its current form) that is independently
defended in contemporary philosophy by writers with no Islamic
agenda. The Islamic position is therefore not a deviation from
contemporary intellectual seriousness; it is a more elaborated
version of a position that contemporary intellectual seriousness has
arrived at.

The reader who is not Muslim is not asked to accept Islamic
theology. They are invited to consider that the Islamic
metaphysical tradition has been articulating, in its own
vocabulary, a position that contemporary Western philosophy has,
through independent argument, found to be the most rigorous response
to the hard problem of consciousness. The two traditions are
intellectual partners, not rivals.

\section{What comes next: Chapter 9 and the Akbarian extension}
\label{sec:idealism-ahead}

Chapter 8 will treat non-computability of consciousness (the
Penrose argument from G\"odel) before we return to the metaphysical
project. Then Chapter 9 will develop the Akbarian metaphysical
tradition, which is the Islamic articulation of the same ontology
that this chapter has presented in contemporary form.

The Akbarian framework extends the analytical-idealist position in
several ways. It articulates the divine names as principles by
which absolute being unfolds into determinate beings (a structural
account of how Mind-at-Large differentiates into the kinds of
mental processes we observe). It articulates \arb{tajalli}
(self-disclosure) as the active mechanism of this unfolding. It
articulates the imaginal world (\arb{`alam al-mithal}) as an
intermediate ontological domain that mediates between pure mental
process and physical appearance. And it articulates the perfect
human (\arb{insan al-kamil}) as the locus that manifests the
underlying unity in its full structure.

These extensions are theological as well as philosophical. They are
not implied by Kastrup's analytical idealism alone. But they are
consistent with it and they enrich the basic analytical-idealist
ontology with substantial philosophical content. Chapter 9 develops
them in detail.

The reader who has worked through this chapter is now prepared for
both the Penrose chapter (Chapter 8) and the Akbarian chapter
(Chapter 9). Each will assume that the basic move from materialism
to idealism is on the table; both will extend the position in
specific directions.

\begin{summarybox}[Chapter summary]
\textbf{The hard problem of consciousness}: a complete physical
description of the brain does not appear to entail the existence of
subjective experience. This problem has resisted solution for
decades.

\textbf{Major positions}: materialism (the dominant academic view,
struggles with the hard problem), property dualism (cleaner but
multiplies fundamental laws), panpsychism (faces the combination
problem), idealism (mind is fundamental).

\textbf{Materialism's difficulties}: the hard problem itself, the
unity of consciousness, mental causation, brain damage anomalies,
psychedelic findings, quantum-mechanical observer effects. Each
admits of materialist response; the cumulative weight is substantial.

\textbf{Analytical idealism (Kastrup)}: there is one fundamental
mind (Mind-at-Large); individual minds are dissociated alters;
physical reality is the extrinsic appearance of mental processes.

\textbf{The dissociation analogy}: clinical dissociative identity
disorder shows that a single mind can appear as multiple
mutually-isolated mental processes. The analytical-idealist
hypothesis extends this to the cosmic scale.

\textbf{Empirical anomalies for materialism, accommodated by
idealism}: hydrocephalic cases, psychedelic findings, near-death
experiences, quantum observer effects.

\textbf{Structural correspondence with Islamic tawhid}: the
ontological monism of Kastrup's Mind-at-Large structurally
corresponds to the Akbarian doctrine of \arb{wahdat al-wujud}. The
two traditions, working from different starting points, have
arrived at structurally similar positions.

\textbf{Differences}: Islamic theology adds substantial content
(personal Allah, revelation, eschatology) that Kastrup's
philosophical argument does not establish. The convergence is in
basic ontology, not full theological content.

\textbf{What comes next}: Chapter 8 treats the Penrose argument
from G\"odel, the fourth major resource of Part II.
\end{summarybox}

\cleardoublepage
